% ============================================================
% Section 101: 理想导体中的电荷弛豫
% ============================================================
\section{电磁学：理想均匀导体中的电荷弛豫}
\label{sec:charge-relaxation}

\begin{modelbox}[题目：导体中空间电荷的时间演化]
一理想均匀导体，其介电常数 $\varepsilon$ 和电导率 $\sigma$ 均为常数。在 $t=0$ 时刻，其中有一电荷分布 $\rho(\bm{r})$，求随后的 $\rho(\bm{r},t)$。
\end{modelbox}

\begin{tipbox}[考点快速分析]
\begin{itemize}
    \item \textbf{电荷守恒}：$\partial\rho/\partial t+\nabla\cdot\bm{j}=0$
    \item \textbf{高斯定律}：$\nabla\cdot\bm{D}=\rho\;\Rightarrow\;\varepsilon\nabla\cdot\bm{E}=\rho$
    \item \textbf{欧姆定律}：$\bm{j}=\sigma\bm{E}$
    \item \textbf{介电弛豫时间}：$\tau_d=\varepsilon/\sigma$
\end{itemize}
\end{tipbox}

\begin{derivationbox}[标准解题过程]
\textbf{Step 1: 建立微分方程}

导体中的电荷守恒方程
\begin{equation}
\frac{\partial\rho}{\partial t}+\nabla\cdot\bm{j}=0.
\end{equation}

由欧姆定律的微分形式 $\bm{j}=\sigma\bm{E}$，得
\begin{equation}
\nabla\cdot\bm{j}=\sigma\nabla\cdot\bm{E}.
\end{equation}

由高斯定律 $\nabla\cdot\bm{D}=\rho$，即 $\varepsilon\nabla\cdot\bm{E}=\rho$，故
\begin{equation}
\nabla\cdot\bm{E}=\frac{\rho}{\varepsilon}.
\end{equation}

代入得
\begin{equation}
\nabla\cdot\bm{j}=\frac{\sigma\rho}{\varepsilon}.
\end{equation}

电荷守恒方程化为
\begin{equation}
\frac{\partial\rho}{\partial t}+\frac{\sigma}{\varepsilon}\rho=0.
\end{equation}

\textbf{Step 2: 求解并得到弛豫时间}

这是一阶线性常微分方程，对每一空间点独立成立。通解为
\begin{equation}
\rho(\bm{r},t)=\rho(\bm{r},0)\,e^{-\sigma t/\varepsilon}=\rho(\bm{r})\,e^{-t/\tau_d},
\end{equation}
其中定义\textbf{介电弛豫时间}（或麦克斯韦弛豫时间）
\begin{equation}
\boxed{\tau_d=\frac{\varepsilon}{\sigma}.}
\end{equation}

\textbf{Step 3: 物理意义与数值估计}

$\tau_d$ 表征导体内部空间电荷自发消散、恢复电中性的特征时间。电导率越高，$\tau_d$ 越短。

实例——半导体锗（Ge）：$\varepsilon_r=15$，$\sigma=0.1\,(\Omega\cdot\text{cm})^{-1}$，
\begin{equation}
\varepsilon=\varepsilon_r\varepsilon_0=15\times8.854\times10^{-14}\approx1.33\times10^{-12}\,\text{F/cm},
\end{equation}
\begin{equation}
\tau_d=\frac{1.33\times10^{-12}}{0.1}\approx1.3\times10^{-11}\,\text{s}.
\end{equation}

任何局域电荷积累在约 $10^{-11}\,$s 内消散。远长于此时间尺度的现象中，半导体内部可视为处处电中性——这正是分析载流子输运和复合过程的\textbf{准电中性条件}的基础。
\end{derivationbox}

\begin{formulabox}[答案汇总]
\begin{equation}
\boxed{\rho(\bm{r},t)=\rho(\bm{r})\,e^{-t/\tau_d},\qquad\tau_d=\frac{\varepsilon}{\sigma}.}
\end{equation}
\end{formulabox}

\begin{insightbox}[物理图像]
\begin{itemize}
    \item 导体中空间电荷的指数衰减是麦克斯韦方程与欧姆定律的直接推论，不依赖于具体几何形状。
    \item 对于良导体（如铜），$\tau_d\sim10^{-19}\,$s 极短；对于弱导电材料（如本征半导体），$\tau_d$ 较长但仍为纳秒至微秒量级。
    \item 介电弛豫概念广泛应用于等离子体物理、半导体器件物理和电磁场理论中。在分析 pn 结、MOS 结构等器件时，准电中性假设是简化计算的关键前提。
\end{itemize}
\end{insightbox}
